\begin{frame}{Natura 2000 og Ramsarområder i forhold til jagt}

\begin{itemize}
  \item \textbf{Natura 2000}: EU’s netværk af beskyttede områder (Habitatsdirektivet + Fuglebeskyttelsesdirektivet). I Danmark: habitatområder og fuglebeskyttelsesområder.
  \item \textbf{Ramsarområder}: Vådområder af international betydning (Ramsar-konventionen). Danmark har ca. 28 områder (ca. 7.300 km²).\\
        Alle danske Ramsarområder indgår også i Natura 2000-netværket.
\end{itemize}

\vspace{0.4cm}

\textbf{Jagt i områderne}
\begin{itemize}
  \item Der er \textbf{ikke} automatisk jagtforbud i Natura 2000- eller Ramsarområder.
  \item Fuglebeskyttelsesdirektivet anerkender jagt som legitim bæredygtig udnyttelse (art. 7 + bilag II).
  \item Jagten skal være bæredygtig og må ikke skade de arter eller naturtyper, områderne er udpeget for.
\end{itemize}

\vspace{0.3cm}

\textbf{Vigtig undtagelse}
\begin{itemize}
  \item Fra 1. juli 2022 er jagt på \textbf{edderfugl} forbudt i de Natura 2000-områder, hvor arten er på udpegningsgrundlaget.
\end{itemize}

\vspace{0.3cm}
\footnotesize
\textit{Øvrige restriktioner kan gælde via vildtreservater, motorbådsjagtregler m.m.}

\end{frame}